\chapter{大模型的决策大脑：Logits、温度系数极限证明与采样算法}

\begin{mathBox}[核心数学定理推导]
\textbf{阅读前须知}：经过了前面四篇的学习，大模型的 Transformer 网络已经完成了前向传播，把前文所有的上下文浓缩进了一个最终的隐藏层向量中。在这最后一公里，这个向量是如何变成真实世界的文字的？为什么温度调高了模型会胡思乱想？调低了就会一本正经？本篇将为你带来严格的数学极限证明与算法解析。
\end{mathBox}

\vspace{0.8em}\hrule\vspace{0.8em}

\section{[目标] 一、最后的终点线：LM Head（语言模型头）}

在模型的最后一层，针对序列的最后一个位置，输出一个经过归一化的隐藏特征向量：

\[
h_{\text{last}} \in \mathbb{R}^{d_{\text{model}}} \quad \text{(例如 } 4096 \text{ 维)}
\]

这个向量就像是一团高度压缩的“思考精华”，里面包含了模型对“下一个字应该是什么”的所有语义判断。
但它还不是人类能看懂的字，我们需要把它映射回\textbf{整个词表（Vocabulary）}。

\subsection{词表矩阵乘法}

假设模型的词表里包含 $V$ 个词（现代大模型如 Llama-3 的词表大小 $V = 128,256$；Qwen2.5 小尺寸 151,936、7B 及以上 152,064）。
我们使用一个投影矩阵 $W_{\text{head}} \in \mathbb{R}^{d_{\text{model}} \times V}$ 与之相乘：

\[
z = h_{\text{last}} W_{\text{head}} \in \mathbb{R}^V
\]

这个输出向量 $z$ 的长度恰好等于词表大小 $V$。
向量里的每一个数值，被称为一个 \textbf{Logit（未归一化对数打分）}：

\[
z = [z_1, z_2, \dots, z_V]
\]

\begin{itemize}
  \item $z_{1024}$ 对应词表里的“的”，数值可能是 $14.5$（极高）；
  \item $z_{305}$ 对应词表里的“烤鸭”，数值可能是 $8.2$；
  \item $z_{9999}$ 对应词表里的“宇宙飞船”，数值可能是 $-3.1$（极低）。
\end{itemize}

现在，模型拿到了对全人类十几万个词汇的打分单。接下来的关键问题是：\textbf{怎么做最终的选词决策？}

\vspace{0.8em}\hrule\vspace{0.8em}

\section{[剖析] 二、决策方法 1：贪婪搜索（Greedy Search）}

最直观的方法，就是直接挑选打分最高的那一个词：

\[
\hat{y} = \arg\max_{i} (z_i)
\]

\begin{center}
\small
\begin{tabularx}{0.7\textwidth}{l c X}
\toprule
\textbf{候选 Token} & \textbf{模型输出 Logit 打分} & \textbf{选择策略判定} \\
\midrule
\texttt{"的"} & $14.5$ & \textbf{最高分！Greedy 策略直接选中} \\
\texttt{"是"} & $12.1$ & 次高分 \\
\texttt{"在"} & $9.3$ & 低分 \\
\texttt{"飞船"} & $-2.0$ & 极低概率 \\
\bottomrule
\end{tabularx}
\end{center}

\subsection{优缺点分析：}

\begin{itemize}
  \item \textbf{优点}：计算最快，结果 100\% 确定，非常适合\textbf{数学推理、编写代码、事实性问答}等有唯一确定答案的场景。
  \item \textbf{缺点}：极其死板，缺乏创造力。更致命的是，贪婪算法只看眼前单步最优，容易陷入局部最优陷阱，导致模型在长文本生成时陷入可怕的\textbf{无脑复读死循环}（“我认为\dots{}\dots{}我认为\dots{}\dots{}我认为\dots{}\dots{}”）。
\end{itemize}

为了打破这种死板，人类引入了\textbf{随机概率采样（Stochastic Sampling）}。

\vspace{0.8em}\hrule\vspace{0.8em}

\section{[温度] 三、【重磅数学推导】：温度系数（Temperature）的极限证明}

为了在采样中可控地调节模型的“严肃度”与“创造力”，科学家在 Softmax 中引入了一个正实数参数 $T$（Temperature，温度系数）：

\[
p_i(T) = \frac{\exp\left(\frac{z_i}{T}\right)}{\sum_{j=1}^V \exp\left(\frac{z_j}{T}\right)}
\]

你一定在各类大模型界面中调节过这个滑动条（通常在 $0.0 \sim 2.0$ 之间）。
\textbf{但你是否知道，在严谨的数学分析中，温度参数在极限状态下到底意味着什么？}

下面我们给出两个经典的\textbf{数学定理与极限推导}。

\vspace{0.8em}\hrule\vspace{0.8em}

\subsection{定理 1：若最大 logit 唯一，当温度 $T \to 0^+$ 时，概率分布严格退化为贪婪搜索（确定性决策）}

\begin{mathBox}[核心数学定理推导]
\textbf{【证明】}：
设 Logits 向量中存在唯一的最大值 $z_{\max} = \max_j (z_j)$。
对任意其他非最大项 $i \neq \max$，均有：
\end{mathBox}

\[
z_i - z_{\max} < 0
\]

\begin{noteBox}[核心导读与背景]
我们将 Softmax 公式分子与分母同时除以 $\exp\left(\frac{z_{\max}}{T}\right)$（利用平移不变性）：
\end{noteBox}

\[
p_i(T) = \frac{\exp\left(\frac{z_i - z_{\max}}{T}\right)}{\sum_{j=1}^V \exp\left(\frac{z_j - z_{\max}}{T}\right)}
\]

\begin{noteBox}[核心导读与背景]
考察当 $T \to 0^+$ 时的极限行为：

\textbf{情况 A：对于最大值项自身（$i = \max$）}：
\end{noteBox}

\[
z_{\max} - z_{\max} = 0 \implies \exp\left(\frac{0}{T}\right) = e^0 = 1
\]

\begin{noteBox}[核心导读与背景]
\textbf{情况 B：对于任意非最大值项（$i \neq \max$）}：
因为 $z_i - z_{\max} < 0$（为一个严格负常数 $-c$），当 $T \to 0^+$ 时：
\end{noteBox}

\[
\lim_{T \to 0^+} \frac{z_i - z_{\max}}{T} = -\infty
\]

\begin{noteBox}[核心导读与背景]
根据自然指数极限：
\end{noteBox}

\[
\lim_{T \to 0^+} \exp\left(\frac{z_i - z_{\max}}{T}\right) = \lim_{x \to -\infty} e^x = \mathbf{0}
\]

\begin{noteBox}[核心导读与背景]
现在把分子和分母的极限合并：
- 分母的总和：
\end{noteBox}

\[
\lim_{T \to 0^+} \sum_{j=1}^V \exp\left(\frac{z_j - z_{\max}}{T}\right) = 1 + \underbrace{0 + 0 + \dots + 0}_{V-1 \text{ 个}} = \mathbf{1}
\]

\begin{noteBox}[核心导读与背景]
- 最大值项的概率：
\end{noteBox}

\[
\lim_{T \to 0^+} p_{\max}(T) = \frac{1}{1} = \mathbf{1.0} \quad (100\%)
\]

\begin{noteBox}[核心导读与背景]
- 所有非最大值项的概率：
\end{noteBox}

\[
\lim_{T \to 0^+} p_{i \neq \max}(T) = \frac{0}{1} = \mathbf{0.0} \quad (0\%)
\]

\begin{noteBox}[核心导读与背景]
\textbf{Q.E.D. 证毕！}
\end{noteBox}

\subsubsection{物理直觉：}

\begin{warningBox}[避坑指南与核心警示]
{}[注意] \textbf{"最大值唯一"这个前提不能省}：如果有 $k$ 个词的 logit 并列最大，同样的推导给出极限是\textbf{这 $k$ 个词各占 $1/k$ 的均匀分布}，而不是确定性地选其中一个。真实模型里 logit 恰好相等很少见，但 BF16 精度下前几名"数值上相等"并不罕见，这也是低温采样有时仍有随机性的原因之一。
\end{warningBox}

\textbf{温度越低（$T \to 0$），最大值项的优势被无限放大，概率瞬间形成尖锐的独峰（全部概率集中在一个词上的"单点分布"，即 one-hot），模型 100\% 必然挑选中打分最高的词！这就是为什么调低温度能让回答严谨、稳定、绝不胡说八道！}

\vspace{0.8em}\hrule\vspace{0.8em}

\subsection{定理 2：当温度 $T \to +\infty$ 时，概率分布严格退化为均匀分布（纯随机乱码）}

\begin{mathBox}[核心数学定理推导]
\textbf{【证明】}：
考察指数内部项在 $T \to +\infty$ 时的极限：
对于任意有限实数 $z_i$：
\end{mathBox}

\[
\lim_{T \to +\infty} \frac{z_i}{T} = 0
\]

\begin{noteBox}[核心导读与背景]
因此：
\end{noteBox}

\[
\lim_{T \to +\infty} \exp\left(\frac{z_i}{T}\right) = e^0 = 1
\]

\begin{mathBox}[核心数学定理推导]
代入 Softmax 公式：
\end{mathBox}

\[
\lim_{T \to +\infty} p_i(T) = \frac{1}{\sum_{j=1}^V 1} = \frac{1}{V}
\]

\begin{noteBox}[核心导读与背景]
\textbf{Q.E.D. 证毕！}
\end{noteBox}

\subsubsection{物理直觉：}

\textbf{温度无穷大（$T \to \infty$）时，词表中每一个词（不论是常见字还是生僻乱码）被选中的概率全部变成了毫无差别的 $1/V$！模型丧失了一切前向推理学习到的知识，完全退化为一个均匀扔骰子的随机数发生器！}

\vspace{0.8em}\hrule\vspace{0.8em}

\section{[目标] 四、工业级截断双子星：Top-K 与 Top-P}

单纯靠 Temperature 调节依然有风险：即使把温度设为正常的 $0.7$，长尾里那些概率只有 $0.0001\%$ 的荒谬生僻字，依然有极其微小的概率被“不幸抽中”，导致模型偶尔冒出一句胡话。

为了彻底扫清这些低质噪声，工业界标配了两种截断算法：

\subsection{Top-K 采样}

\begin{itemize}
  \item \textbf{逻辑}：不管词表有多大，\textbf{永远只保留概率最高的前 $K$ 个候选词}（例如 $K=40$ 或 $K=50$）。
  \item 将第 $K+1$ 到第 $V$ 个词的 Logit 强行赋值为 $-\infty$（概率归零）。
  \item 在保留的这 $K$ 个词中重新归一化 Softmax。
\end{itemize}

\begin{center}
\small
\begin{tabularx}{\textwidth}{X c X}
\toprule
\textbf{高概率候选词（保留采样池）} & \textbf{截断阈值} & \textbf{长尾低概率词（强制置 $-\infty$）} \\
\midrule
前 $K$ 个最高概率词: \texttt{[词 1, 词 2, $\dots$, 词 $K$]} & $\longleftrightarrow$ & 剩余词: \texttt{[词 $K+1$, $\dots$, 词 $V$]} 全部过滤 \\
\bottomrule
\end{tabularx}
\end{center}

\vspace{0.8em}\hrule\vspace{0.8em}

\subsection{Top-P 采样（核采样 Nucleus Sampling，更智能的动态截断）}

Top-K 有一个天然的死穴：\textbf{固定 $K$ 值无法适应不同语境的灵活性！}
\begin{itemize}
  \item \textbf{场景 A}：用户输入“人工智能的英文缩写是\_\_\_\_”。
  \item 这种问题几乎只有唯一正确答案“AI”。此时理想的候选池大小应该是 $K=1$。如果固定 $K=50$，模型反而可能把后 49 个错误答案放进抽奖箱。
  \item \textbf{场景 B}：用户输入“请描述一个美丽的黄昏景色：\_\_\_\_”。
  \item 此时有成百上千个优美的辞藻（“霞光”、“落日”、“晚风”、“云彩”）都是合理的，此时 $K=50$ 又显得太少，扼杀了多样性。
\end{itemize}

\textbf{Top-P 的破局之道（Holtzman et al., ICLR 2020）}：
\begin{noteBox}[核心导读与背景]
\textbf{不要固定候选词的个数，而是固定候选词的「累积置信度阈值 $P$」（例如 $P = 0.9$ 或 $90\%$）！}
\end{noteBox}

\begin{center}
\small
\begin{tabularx}{\textwidth}{p{3.5cm} c X}
\toprule
\textbf{语境场景} & \textbf{候选 Token 概率累积} & \textbf{动态候选池判定} \\
\midrule
\textbf{极确定语境} (如代码/事实) & \texttt{"AI"}: 92\% $\ge 90\%$ & \textbf{立即截断！候选池仅保留 1 个词} \\
\midrule
\textbf{发散创意语境} (如诗歌/创作) & \texttt{"落日"}: 25\%, \texttt{"晚霞"}: 20\%, $\dots$ & \textbf{累积至 90.5\% 截断！动态保留 15 个候选词} \\
\bottomrule
\end{tabularx}
\end{center}

\begin{itemize}
  \item \textbf{语境确定时}：头部词概率极高，累积迅速破 90\%，候选池自动收缩为 1\textasciitilde{}2 个词，极其精准；
  \item \textbf{语境发散时}：头部词概率分散，需要累积几十个词才破 90\%，候选池自动扩大，充分保障多样性！
\end{itemize}

\vspace{0.8em}\hrule\vspace{0.8em}

\section{五、防复读机制：重复惩罚（Repetition Penalty）}

你一定见过模型陷入“复读机”状态。怎么在数学上给它一记当头棒喝？

在生成第 $t$ 步时，算法会检查前文中\textbf{已经出现过的所有词}的索引集合 $\mathcal{S}_{\text{history}}$。
对于集合中的词，对其打分 $z_i$ 施加一个惩罚系数 $\theta > 1.0$（通常取 $1.1 \sim 1.3$）：

\[
z_i' = \begin{cases}
z_i / \theta & \text{若 } z_i > 0 \quad (\text{主动压低正向打分}) \\
z_i \times \theta & \text{若 } z_i \le 0 \quad (\text{加大负向惩罚力度})
\end{cases}
\]

\textbf{通过这一轻巧的代数缩放，曾经说过的词在下一步重新争夺第一名的难度被大幅提高，模型被迫转向寻找新的表达方式，显著缓解复读机问题。}

\begin{warningBox}[避坑指南与核心警示]
{}[注意] 但它是一把钝器：它惩罚的是"出现过的所有 token"，不区分该不该重复。人名、专有名词、标点、代码里的变量名本来就需要反复出现，惩罚过重会让模型改用同义词、写错变量名，甚至破坏格式。而且它连提示词里的 token 也一起惩罚（Lab 07b 实验 5 里 HF 的实现就是这样）。所以实际服务中惩罚系数通常只取 1.05\textasciitilde{}1.2，另有 frequency / presence penalty 等更细的变体。
\end{warningBox}

\vspace{0.8em}\hrule\vspace{0.8em}

\section{[检验] 前七篇小结}

到这里，你已经\textbf{从零起点、不依赖任何黑盒}，贯通了"推理在算什么"的整条链路：

\begin{enumerate}
  \item 第 1 篇：计算机如何从规则、统计到被 Transformer 彻底征服
  \item 第 2 篇：向量空间几何、点积的几何意义与 Safe Softmax 溢出防护
  \item 第 3 篇：导数、反向传播、交叉熵，Softmax 雅可比与 $p - y$ 梯度
  \item 第 4 篇：严格证明 Attention 为什么必须除以 $\sqrt{d_k}$，推导 RoPE 相对位置内积恒等性
  \item 第 5 篇：从原版 Transformer 到 Llama / DeepSeek，每一处改动的历史动机
  \item 第 6 篇：推理与训练的本质区别，代数展开 KV Cache 消除冗余并剖析显存挑战
  \item 第 7 篇：温度参数在 $T \to 0$ 和 $T \to \infty$ 下的数学极限，Top-K、Top-P 与重复惩罚
\end{enumerate}

但要真正理解推理\textbf{为什么快、为什么慢、怎么优化}，还需要继续：
{}\textbf{第 8 篇：硬件视角} $\rightarrow$ \textbf{第 9 篇：FlashAttention} $\rightarrow$ \textbf{第 10 篇：投机解码} $\rightarrow$ \textbf{第 11 篇：量化} $\rightarrow$ \textbf{第 12 篇：推理服务系统} $\rightarrow$ \textbf{第 13 篇：多卡并行}

完整顺序（含实验与外部资源）见 \textbf{LEARNING\_PATH.md}。

\vspace{0.8em}\hrule\vspace{0.8em}

\textbf{上一篇}：\textbf{第 6 篇：什么是推理？自回归生成全生命周期与 KV Cache 代数证明} ｜ \textbf{下一篇}：\textbf{第 8 篇：硬件视角：GPU 存储层级、FLOPs、算术强度与浮点格式} ｜ \textbf{总路线}：\textbf{LEARNING\_PATH.md}
